\documentclass[11pt]{amsart}
\usepackage[margin=1.05in]{geometry}
\usepackage{amsmath,amssymb,amsthm,mathtools}
\usepackage{enumitem}
\usepackage[hidelinks]{hyperref}

\newtheorem{theorem}{Theorem}[section]
\newtheorem{lemma}[theorem]{Lemma}
\newtheorem{proposition}[theorem]{Proposition}
\newtheorem{corollary}[theorem]{Corollary}
\theoremstyle{definition}
\newtheorem{definition}[theorem]{Definition}
\newtheorem*{remark}{Remark}
\newtheorem{question}{Question}
\newtheorem*{thmR}{Theorem R}

\newcommand{\R}{\mathbb{R}}
\newcommand{\C}{\mathbb{C}}
\newcommand{\Vint}{V_{\mathrm{int}}}
\newcommand{\Ran}{\operatorname{Ran}}
\newcommand{\Gram}{\operatorname{Gram}}
\newcommand{\norm}[1]{\lVert #1\rVert}
\newcommand{\abs}[1]{\lvert #1\rvert}
\newcommand{\ip}[2]{\langle #1,#2\rangle}
\newcommand{\PMTA}{\textsf{PMT-A}}

\title[Theorem R]{Theorem R: Sharp Topology-Independent Stability of Real Projected Multilinear Trees}
\author{Review note}
\date{Version v2 (tag \texttt{theorem-R-v2-review}, commit \texttt{3d920f9}), 13 September 2026.
Status: advisory proof draft; author audit only; not independently reviewed; novelty not established.}

\begin{document}

\begin{abstract}
A projected multilinear tree composes independently chosen multilinear maps
between finite-dimensional real Hilbert spaces along a rooted tree and projects
orthogonally after every internal operation. Under a uniform operator-norm budget
$M$ and a closure budget $\rho=\eta M$ on the part of each law that leaves the
retained subspace on projected inputs, let $C^P_T(\eta)$ be the sharp constant for
the projected output error. We prove that
\[
C^P_T(\eta)=\frac1\eta\max_{0\le\theta\le\arcsin\eta}\bigl|1-(\cos\theta\,e^{i\theta})^{k-1}\bigr|,
\]
where $k$ is the number of internal nodes of $T$. The constant depends on the number
of internal nodes and on the leakage, but not on depth or topology. The upper bound
rests on a Löwner-order state invariant and a tensor angle inequality for projective
norms; the lower bound is an explicit construction valid for every tree.
\end{abstract}

\maketitle

\section*{Introduction}

\subsection*{The problem}
A computation that alternates multilinear operations with orthogonal reductions is
modelled by a rooted tree. Each internal node applies a multilinear law to its
children's values, and a projector keeps only a retained subspace. Examples include
truncated tensor-network contractions, Galerkin projections of polynomial
nonlinearities, and reduced models of recursive compositions. The recursively
reduced value $R_r$ is compared with the ambient value $F_r$ computed without any
reduction. The natural question is extremal: if every law has norm at most $M$ and
leaks at most $\rho$ out of its retained subspace on already-reduced inputs, how
large can the projected error $\norm{P_rF_r-R_r}$ be?

\subsection*{What was known}
Telescoping over arguments gives the universal bound
\begin{equation}\label{eq:universal}
\norm{P_rF_r-R_r}\le(k-1)\,\rho\,M^{k-1}\prod_\ell\norm{z_\ell}.
\end{equation}
Indeed, by induction $\norm{F_v-R_v}\le k_v\rho M^{k_v-1}L_v$ at every node. The
difference $\mu_v(F_c)-\mu_v(R_c)$ telescopes into one term per child, and the local
defect $Q_v\mu_v(R_c)$ contributes one more $\rho$; at the root this defect is
annihilated by $P_r$. The bound~\eqref{eq:universal} is asymptotically sharp as
$\eta=\rho/M\to0$. For $k=2$ it is exact ($C_2^P\equiv1$). For $k=3$ the exact
constant is $W_3(\eta)=\sqrt{4-3\eta^2}$ for $\eta\le\sqrt{2/3}$ and $2/(\sqrt3\eta)$
above, for both three-node skeletons. This shows that orthogonality forbids
simultaneous saturation of all local errors at positive leakage.

\subsection*{What is proved here}
For every tree the sharp constant is a single explicit function of the number $k$ of
internal nodes and the leakage $\eta$ (Theorem R). The whole family of trees collapses
to the one-parameter sequence $C_1,C_2,C_3,\dots$. Section~\ref{sec:consequences}
records the resulting closed forms, the second-order coefficient, and the behaviour
as $k\to\infty$.

\subsection*{The idea}
The state of a node is the $2\times2$ Gram matrix of its ambient and reduced values.
States are ordered by Löwner domination against the Gram matrix of a unit vector and
a vector of norm $\rho$ at an angle at most $\chi$ (order O2). Two facts make this
order propagate through the tree:
\begin{itemize}
\item a projection step multiplies the chain label $\rho e^{i\chi}$ by
  $w(\theta)=\cos\theta\,e^{i\theta}$, because isometric dilations preserve angles and
  the closure budget;
\item a multilinear node multiplies the labels of its children, by a tensor angle
  inequality for projective norms.
\end{itemize}
The labels therefore compose as complex numbers, and the error is at most
$\abs{1-\prod w(\theta_v)}$ over the $k-1$ non-root internal nodes. A concavity
argument then shows that equal angles are extremal. Conversely, multiplication in
$\R^2\cong\C$ realizes every such product exactly, for every tree.

\subsection*{Relation to the literature}
The triangle inequality for angles under nonexpansive maps also underlies the scaled
relative graph calculus~\cite{RHY}. Isometric dilations of contractions are
classical~\cite{NF}, and projective tensor norms are standard~\cite{Ryan}. We are not
aware of the tensor angle inequality (Theorem~\ref{thm:TA}) or of Theorem R in the
literature, but a systematic comparison has not been carried out, and no novelty is
claimed.

\subsection*{How to read this note}

The argument splits into two independent halves:
\begin{align*}
\text{upper bound: }&\text{Definitions}\to\text{O2}\to\text{Tensor angle lemma}\to\text{Lemma 3}\\
&\to\text{Multiplicative envelope}\to\text{Diagonal Lemma},\\
\text{lower bound: }&\text{Universal sharp witness}.
\end{align*}
Sections~\ref{sec:tensor} and~\ref{sec:diagonal} use no tree vocabulary and can be
read on their own. The step we regard as the one most in need of independent
scrutiny is the passage \emph{tensor angle inequality $\Rightarrow$ preservation
of the order O2 at a multilinear node} (Theorem~\ref{thm:TA} and
Lemma~\ref{lem:lemma3}). It is also available as a separate short note.

Throughout, $k$ denotes the \emph{number of internal nodes}, not the depth. A chain
with $k$ internal nodes has depth $k$; a star with $k$ internal nodes has depth $2$.

\section{Definitions}\label{sec:defs}

\subsection{The class \PMTA}

\begin{enumerate}[label=\textup{(A\arabic*)}]
\item $T$ is a finite rooted tree with internal nodes $\Vint$, $k:=\abs{\Vint}\ge1$,
  root $r$, and leaves $L$. Each internal node $v$ has $m_v\ge1$ ordered children,
  some of which may be leaves and some internal.
\item Every node $w$ carries a finite-dimensional real Hilbert space $H_w$ of
  arbitrary dimension (dimension $2$ is allowed).
\item Every internal $v$ carries an orthogonal projector $P_v$ of arbitrary rank.
  Rank $0$, rank $1$ and $P_v=I$ are all allowed, also at the root. Write $Q_v:=I-P_v$.
\item Every internal $v$ carries a multilinear map $\mu_v$ from the product of its
  children's spaces to $H_v$, with
  $\norm{\mu_v}:=\sup\{\norm{\mu_v(x_1,\dots)}:\norm{x_i}\le1\}\le M$.
  Laws at different nodes are chosen independently.
\item Leaves carry nonzero data $z_\ell$. Set $F_\ell=R_\ell=z_\ell$, and for internal $v$
  \[
  F_v=\mu_v(F_{\mathrm{children}}),\qquad R_v=P_v\,\mu_v(R_{\mathrm{children}}).
  \]
\item \emph{Closure.} For every internal $v$,
  $\norm{Q_v\mu_v(x_1,\dots)}\le\rho\prod\norm{x_i}$ whenever each internal-child
  slot receives a vector in that child's $\Ran P$. Leaf slots receive arbitrary
  vectors. Here $0<\rho\le M$ and $\eta:=\rho/M\in(0,1]$; the value $\eta=0$ is not
  part of the class.
\item The projected error is $E^P_T:=\norm{P_rF_r-R_r}$, and
  \[
  C^P_T(\eta):=\sup\frac{E^P_T}{\rho\,M^{k-1}\prod_\ell\norm{z_\ell}},
  \]
  the supremum over all admissible realizations, dimensions and ranks.
\end{enumerate}

\subsection{Normalization and trajectory closure}

\begin{enumerate}[label=\textup{(N\arabic*)}]
\item \emph{Scaling.} The ratio is invariant under $\mu_v\mapsto\lambda\mu_v$ and
  $z_\ell\mapsto\nu_\ell z_\ell$. We take $M=1$, $\rho=\eta$ and unit leaves.
\item \emph{Leaf freezing.} Fix leaf arguments at the leaf data. At each internal $v$
  with internal children $c_1,\dots,c_m$ ($m\ge0$) this gives a multilinear map on
  the internal slots with norm $\le1$, the same $F$ and $R$, and
  \[
  \textup{(TC)}\qquad \norm{Q_v\mu_v(R_{c_1},\dots,R_{c_m})}\le\eta\prod_i\norm{R_{c_i}},
  \]
  because $R_{c_i}\in\Ran P_{c_i}$. For $m=0$, $\mu_v$ is a vector $f_v$ with
  $\norm{f_v}\le1$ and $\norm{Q_vf_v}\le\eta$.
\end{enumerate}

The upper bound uses only norms $\le1$, orthogonality of each $P_v$, and (TC). The
witness satisfies the full closure (A6).

\subsection{The theorem}

Let
\[
w(\theta):=\cos\theta\,e^{i\theta}=\tfrac12\bigl(1+e^{2i\theta}\bigr),\qquad
C_k(\eta):=\frac1\eta\max_{0\le\theta\le\arcsin\eta}\bigl|1-w(\theta)^{k-1}\bigr|.
\]

\begin{thmR}
For every finite rooted tree $T$ in real \PMTA{} with $k$ internal nodes and every
$\eta\in(0,1]$,
\[
C^P_T(\eta)=C_k(\eta).
\]
\end{thmR}

\begin{remark}
At $k=2,3$ this recovers $C_2^P=1$ and the known two-regime constant
$W_3(\eta)=\sqrt{4-3\eta^2}$ for $\eta\le\sqrt{2/3}$, $2/(\sqrt3\eta)$ above.
As $\eta\downarrow0$, $C_k(\eta)=(k-1)-\frac{(k-1)(k-2)(k+6)}{24}\eta^2+O(\eta^4)$.
\end{remark}

\section{The order O2}\label{sec:O2}

For $\rho\ge0$ and $\psi\in[0,\pi]$ put
\[
K(\rho,\psi):=\begin{pmatrix}1&\rho\cos\psi\\ \rho\cos\psi&\rho^2\end{pmatrix},
\]
the Gram matrix of a unit vector and a vector of norm $\rho$ at angle $\psi$.
Angles are $\angle(a,b)=\arccos\bigl(\ip ab/(\norm a\norm b)\bigr)\in[0,\pi]$; this is a
metric on the unit sphere. $\preceq$ denotes the Löwner order.

\begin{definition}[O2]\label{def:O2}
A pair $(F,R)$ is \emph{dominated by the label} $(\rho,\chi)$, with $\rho,\chi\ge0$, if
\[
\Gram(F,R)\preceq K(\rho,\psi')\quad\text{for some }\psi'\in[0,\min(\chi,\pi)].
\]
The label $\chi$ is a cumulative phase bookkeeping parameter. It is not
reconstructed from the Gram matrix.
\end{definition}

\begin{lemma}[Realization]\label{lem:realization}
Assume $\Gram(F,R)\preceq K(\rho,\psi)$ with $\rho\ge0$, $\psi\in[0,\pi]$. Let $X$ be a
$2$-dimensional real space with unit vectors $\hat F,\hat R$ at angle $\psi$. Then
there is $B:X\to H$ with $\norm B\le1$, $B\hat F=F$ and $B(\rho\hat R)=R$. In particular
$\norm F\le1$ and $\norm R\le\rho$.
\end{lemma}

\begin{proof}
For real $\alpha,\beta$ we have
$\norm{\alpha F+\beta R}^2\le\norm{\alpha\hat F+\beta\rho\hat R}^2$. So
$B(\alpha\hat F+\beta\rho\hat R):=\alpha F+\beta R$ is well defined and contractive
on the span. Extend it by $0$ on the orthogonal complement. This covers $\rho=0$
(forcing $R=0$) and $\psi\in\{0,\pi\}$ (forcing $R=\pm\rho F$).
\end{proof}

\begin{lemma}[(N$-$) implies O2]\label{lem:N-}
Let $S\in[0,\pi]$ and let $f,g$ satisfy
\[
\textup{(N$-$)}\qquad \norm{af+bg}\le\abs{ae^{iS}+b}\quad\text{for all real }a,b\text{ with }ab\le0.
\]
Then $\Gram(f,g)\preceq K(1,\psi)$ for some $\psi\in[0,S]$.
\end{lemma}

\begin{proof}
Let $G=\Gram(f,g)$.
\begin{enumerate}[label=\textup{(\roman*)}]
\item The choices $b=0$ and $a=0$ give $G_{11},G_{22}\le1$. Put $\mu:=\sqrt{(1-G_{11})(1-G_{22})}$.
\item (N$-$) reads $a^2(1-G_{11})+b^2(1-G_{22})-2\abs{ab}(\cos S-G_{12})\ge0$.
  Minimizing over the ratio $\abs a:\abs b$, this holds for all $a,b$ if and only if
  $G_{12}+\mu\ge\cos S$.
\item Since its diagonal is nonnegative, $K(1,\psi)-G\succeq0$ if and only if
  $\cos\psi\in[G_{12}-\mu,G_{12}+\mu]$.
\item This interval meets $[\cos S,1]$: $G_{12}-\mu\le\sqrt{G_{11}G_{22}}\le1$, and
  $\cos S\le G_{12}+\mu$ by (ii).
\end{enumerate}
Taking $\cos\psi$ in the intersection gives $\psi\in[0,S]$.
\end{proof}

\begin{remark}
(N$-$) bounds the alignment $G_{12}$ only from below. Excess alignment is absorbed by
a dilation angle $\psi<S$. This is the precise sense in which extra alignment
cannot help the adversary.
\end{remark}

\section{The tensor angle lemma}\label{sec:tensor}

For real Hilbert spaces $H_1,\dots,H_m$, $\norm{\cdot}_\pi$ denotes the projective
tensor norm on $H_1\otimes\cdots\otimes H_m$. It is the dual of the multilinear
operator norm: $\norm{\mathsf T}_\pi=\sup W(\mathsf T)$ over $m$-linear forms $W$ with
$\abs{W(x_1,\dots,x_m)}\le\prod\norm{x_i}$.

\begin{theorem}[Tensor angle inequality]\label{thm:TA}
Let $u_i,v_i\in H_i$ be unit vectors, $\psi_i:=\angle(u_i,v_i)$ and
$S:=\sum_{i=1}^m\psi_i$. For all real $a,b$ with $ab\le0$,
\[
\Bigl\|a\bigotimes_{i=1}^mu_i+b\bigotimes_{i=1}^mv_i\Bigr\|_\pi\le\bigl|ae^{i\min(S,\pi)}+b\bigr|.
\]
\end{theorem}

\begin{proof}
\emph{Reductions.} We may take $a=1$, $b=-t$ with $t\ge0$. If $S\ge\pi$, the right side
is $1+t$ and the claim is the triangle inequality, so assume $S<\pi$. Factors of
dimension one may be embedded isometrically into $\R^2$, since projective norms of
Hilbert factors are preserved under isometric embeddings.

\emph{Induction on $m$.} For $m=1$ the claim is the identity
$\norm{u-tv}^2=1+t^2-2t\cos\psi_1$. Let $m\ge2$.
\begin{enumerate}[label=\textup{(\arabic*)}]
\item Put $X=\bigotimes_{i\ge2}u_i$, $Y=\bigotimes_{i\ge2}v_i$ and $S'=S-\psi_1$.
  By associativity of $\otimes_\pi$ and duality, the left side equals
  $\sup_\beta[\beta(u_1,X)-t\beta(v_1,Y)]$ over bilinear forms $\beta$ on
  $H_1\times(\otimes_{i\ge2}H_i)$ with $\abs{\beta(z,W)}\le\norm z\norm W_\pi$.
\item Fix a $2$-plane $\Pi\ni u_1,v_1$. By the Riesz representation,
  $\beta(z,W)=\ip z{\Lambda W}$ for $z\in\Pi$, with $\norm{\Lambda W}\le\norm W_\pi$.
\item Put $p=\Lambda X$ and $q=\Lambda Y$. The induction hypothesis gives
  $\norm{\alpha p+\beta'q}\le\abs{\alpha e^{iS'}+\beta'}$ for $\alpha\beta'\le0$.
\item Lemmas~\ref{lem:N-} and~\ref{lem:realization} give unit vectors $\hat p,\hat q$
  at an angle $\tilde\psi\le S'$ and a contraction $C$ with $C\hat p=p$, $C\hat q=q$.
\item Let $\mathrm{Rot}$ be the rotation of $\Pi$ taking $v_1$ to $u_1$. Then
  $\beta(u_1,X)-t\beta(v_1,Y)=\ip{u_1}{C\hat p-t\,\mathrm{Rot}\,C\hat q}\le\norm{C\hat p-t\,\mathrm{Rot}\,C\hat q}$.
\item Let $V=(C,(I-C^*C)^{1/2})$, an isometry into $\Pi\oplus D$, and
  $\tilde R=\mathrm{Rot}\oplus I$. For unit $y$,
  $\ip y{\tilde Ry}=\norm{y_\Pi}^2\cos\psi_1+\norm{y_D}^2\ge\cos\psi_1$, so
  $\angle(y,\tilde Ry)\le\psi_1$.
\item The $\Pi$-component of $V\hat p-t\tilde RV\hat q$ is
  $C\hat p-t\,\mathrm{Rot}\,C\hat q$. The vectors $V\hat p$ and $t\tilde RV\hat q$ have
  norms $1$ and $t$, and the angle between them is at most $\tilde\psi+\psi_1\le S<\pi$.
\end{enumerate}
Therefore the quantity in (5) is at most $\abs{1-te^{iS}}$.
\end{proof}

\begin{proposition}[Sharpness]\label{prop:TAsharp}
If $S\le\pi$, equality holds in Theorem~\ref{thm:TA}.
\end{proposition}

\begin{proof}
Identify each $\Pi_i\cong\C$ with $v_i\mapsto1$ and $u_i\mapsto e^{i\psi_i}$, and let
$\pi_i$ be the orthogonal projection onto $\Pi_i$. For $\abs\lambda=1$ the form
$W(x)=\mathrm{Re}\bigl(\lambda\prod_i\pi_ix_i\bigr)$ has norm at most $1$, and
$W(\otimes u_i)-tW(\otimes v_i)=\mathrm{Re}\bigl(\lambda(e^{iS}-t)\bigr)$. Choose $\lambda$.
\end{proof}

\section{Lemma 3: O2 is preserved at non-root internal nodes}\label{sec:lemma3}

\begin{lemma}[Child substitution and closure normalization]\label{lem:substitution}
Let $v$ be internal with internal children $c_1,\dots,c_m$, $m\ge1$, let
$\norm{\mu_v}\le1$, let (TC) hold, and let each $(F_{c_i},R_{c_i})$ be dominated by
$(\rho_i,\chi_i)$ via $\psi_i\in[0,\min(\chi_i,\pi)]$. Let $B_i$ be as in
Lemma~\ref{lem:realization}, with unit $\hat F_i,\hat R_i$ at angle $\psi_i$. Put
$\mu'(x):=\mu_v(B_1x_1,\dots,B_mx_m)$ and $\rho:=\prod_i\rho_i$. Then:
\begin{enumerate}[label=\textup{(\roman*)}]
\item $\norm{\mu'}\le1$;
\item $F_v=\mu'(\hat F_1,\dots)$ and $\mu_v(R_{c_1},\dots)=\rho\,\mu'(\hat R_1,\dots)$;
\item if $\rho>0$, then $\norm{Q_v\mu'(\hat R_1,\dots)}\le\eta$;
\item if $\rho=0$, then $\mu_v(R_{c_1},\dots)=0$.
\end{enumerate}
\end{lemma}

\begin{proof}
(i) Composition with contractions. (ii) Multilinearity. (iii) By (ii) and (TC),
$\norm{Q_v\mu'(\hat R)}=\norm{Q_v\mu_v(R_c)}/\rho\le\eta\prod\norm{R_{c_i}}/\rho\le\eta$,
using $\norm{R_{c_i}}\le\rho_i$. (iv) Some $R_{c_i}=0$.
\end{proof}

\begin{lemma}[Lemma 3]\label{lem:lemma3}
Let $v\ne r$ be internal with $m\ge0$ internal children satisfying the hypotheses of
Lemma~\ref{lem:substitution}. Then there is $\theta_v\in[0,\arcsin\eta]$ such that
$(F_v,R_v)$ is dominated by
\[
(\rho_v,\chi_v):=\Bigl(\cos\theta_v\prod_i\rho_i,\;\theta_v+\sum_i\chi_i\Bigr)
\]
(empty product $1$, empty sum $0$). When $\prod_i\rho_i>0$ one can take
\[
\sin\theta_v=\frac{\norm{Q_v\mu_v(R_{c_1},\dots,R_{c_m})}}{\prod_i\rho_i}.
\]
\end{lemma}

\begin{proof}
\emph{Case $\prod_i\rho_i=0$} (so $m\ge1$). By Lemma~\ref{lem:substitution}(iv),
$R_v=0$, and $\norm{F_v}\le1$. Hence $\Gram(F_v,0)\preceq K(0,0)$; take $\theta_v=0$.

\emph{Case $\rho:=\prod_i\rho_i>0$.}
\begin{enumerate}[label=\textup{(\arabic*)}]
\item \emph{Raw output.} Put $f=\mu'(\hat F)$ and $g=\mu'(\hat R)$, so $F_v=f$ and
  $R_v=\rho P_vg$. For $m=0$ put $f=g=f_v$, $\rho=1$, $S=0$; otherwise $S=\sum\psi_i$.
\item \emph{(N$-$) for the raw output.} With $S'=\min(S,\pi)$, for $ab\le0$,
  \[
  \norm{af+bg}=\sup_{\norm\omega\le1}\ip\omega{\mu'(a\otimes\hat F+b\otimes\hat R)}\le\norm{a\otimes\hat F+b\otimes\hat R}_\pi\le\abs{ae^{iS'}+b}
  \]
  by Theorem~\ref{thm:TA}. For $m=0$ this is immediate.
\item \emph{Domination.} Lemmas~\ref{lem:N-} and~\ref{lem:realization} give unit
  $\hat F',\hat E'$ at an angle $\tilde\psi\le S'$ and a contraction $C$ with
  $C\hat F'=f$, $C\hat E'=g$.
\item \emph{Dilation.} Let $V=(C,(I-C^*C)^{1/2})$ and $P'=P_v\oplus I$, and set
  $\tilde F=V\hat F'$, $\tilde R=\rho P'V\hat E'$. Their $H_v$-components are $F_v$
  and $R_v$, so $\Gram(F_v,R_v)\preceq\Gram(\tilde F,\tilde R)$.
\item \emph{Norms.} $\norm{\tilde F}=1$. Also
  $\norm{(I-P')V\hat E'}=\norm{Q_vg}\le\eta$ by Lemma~\ref{lem:substitution}(iii)
  (by (N2) if $m=0$). Put $\theta_v:=\arcsin\norm{Q_vg}$. Then
  $\norm{\tilde R}=\rho\cos\theta_v$, and $\norm{Q_vg}=\norm{Q_v\mu_v(R_c)}/\rho$
  gives the stated formula.
\item \emph{Angle.} If $\tilde R=0$, the pair is dominated by $K(0,0)$. Otherwise
  $\angle(\tilde F,\tilde R)\le\angle(V\hat F',V\hat E')+\angle(V\hat E',P'V\hat E')=\tilde\psi+\theta_v\le\chi_v$,
  so $\Gram(\tilde F,\tilde R)=K(\rho_v,\angle(\tilde F,\tilde R))$.
\end{enumerate}
Combining (4) and (6) proves the claim.
\end{proof}

\section{The multiplicative envelope}\label{sec:envelope}

Encode a label as $z=\rho e^{i\chi}$. Lemma~\ref{lem:lemma3} then reads
$z_v=w(\theta_v)\prod_iz_{c_i}$.

\begin{proposition}[Unrolling]\label{prop:unroll}
Applying Lemma~\ref{lem:lemma3} from the leaves upward, every non-root internal $v$
has $(F_v,R_v)$ dominated by $z_v=\prod_{u\in\Vint(T_v)}w(\theta_u)$, where $T_v$ is the
subtree of $v$. Since the subtrees of the root's internal children are pairwise
disjoint and exhaust $\Vint\setminus\{r\}$,
\[
\prod_iz_{c_i}=\prod_{u\in\Vint\setminus\{r\}}w(\theta_u),
\]
a product with exactly $k-1$ factors.
\end{proposition}

\begin{remark}
The count uses that $T$ is a tree; for directed acyclic graphs this step fails as
stated. O2 induces a commutative multiplicative envelope for tree states. We make
no stronger algebraic claim.
\end{remark}

\begin{lemma}[Root reading]\label{lem:root}
Let the root have $m\ge1$ internal children dominated by $(\rho_i,\chi_i)$ via
angles $\psi_i$, and let $\norm{\mu_r}\le1$. No closure is needed at the root. Then
$E^P_T\le\abs{1-\rho e^{i\min(S,\pi)}}$ with $\rho=\prod\rho_i$ and $S=\sum\psi_i$.
\end{lemma}

\begin{proof}
$E^P_T=\norm{P_r(\mu_r(F_c)-\mu_r(R_c))}$. Substitute the children
(Lemma~\ref{lem:substitution}(i)--(ii)), scalarize with a unit $\omega\in\Ran P_r$, and
apply Theorem~\ref{thm:TA} with $a=1$, $b=-\rho$. If $\rho=0$ the bound is
$1\ge\norm{F_r}$.
\end{proof}

\begin{corollary}\label{cor:envelope}
There are $\theta_u\in[0,\arcsin\eta]$, $u\in\Vint\setminus\{r\}$, with
\[
E^P_T\le\bigl|1-R\,e^{i\min(\Theta,\pi)}\bigr|,\qquad R=\prod_{u\ne r}\cos\theta_u,\quad\Theta=\sum_{u\ne r}\theta_u.
\]
\end{corollary}

\begin{proof}
Lemma~\ref{lem:root} with $S\le\Theta$, since $\abs{1-Re^{i\varphi}}$ is
nondecreasing in $\varphi\in[0,\pi]$ for $R\ge0$.
\end{proof}

\section{The Diagonal Lemma}\label{sec:diagonal}

Throughout this section $n\ge1$ and $\alpha\in(0,\pi/2]$.

\begin{lemma}[Subadditivity of cosine products]\label{lem:subadd}
If $\theta_j\ge0$ and $\sum_j\theta_j\le\pi/2$, then
$\prod_j\cos\theta_j\ge\cos\bigl(\sum_j\theta_j\bigr)$.
\end{lemma}

\begin{proof}
Induction on $n$. All partial sums lie in $[0,\pi/2]$, so the sines and cosines
involved are nonnegative. Then $\cos(a+b)=\cos a\cos b-\sin a\sin b\le\cos a\cos b$,
and multiplying the inductive inequality by $\cos\theta_n\ge0$ preserves it.
\end{proof}

\begin{remark}
The hypothesis $\sum\theta_j\le\pi/2$ matters for large angle sums. With $n=4$ and
$\theta_j=\pi/2-\varepsilon$ we get $\prod\cos\theta_j=\sin^4\varepsilon\approx0$ while
$\cos\sum\theta_j\approx1$. For $\sum\theta_j\in(\pi/2,\pi]$ the inequality holds
trivially, because $\cos\sum\theta_j\le0$.
\end{remark}

\begin{lemma}[Jensen for $\log\cos$]\label{lem:jensen}
$\log\cos$ is strictly concave on $[0,\pi/2)$, since $(\log\cos)''=-\sec^2<0$. Hence
for $\theta_j\in[0,\pi/2)$ with $\Theta=\sum\theta_j$ we have
$\prod_j\cos\theta_j\le\cos^n(\Theta/n)$, with equality if and only if all
$\theta_j$ are equal.
\end{lemma}

\begin{lemma}[Diagonal Lemma]\label{lem:diagonal}
\[
\max_{\theta\in[0,\alpha]^n}\Bigl|1-\prod_{j=1}^nw(\theta_j)\Bigr|=\max_{t\in[0,\alpha]}\bigl|1-w(t)^n\bigr|.
\]
Moreover, for every $\theta\in[0,\alpha]^n$, with $R=\prod\cos\theta_j$ and
$\Theta=\sum\theta_j$,
\[
\bigl|1-Re^{i\min(\Theta,\pi)}\bigr|\le\max_{t\in[0,\alpha]}\abs{1-w(t)^n}.
\]
\end{lemma}

\begin{proof}
The inequality ``$\ge$'' is the diagonal choice. For ``$\le$'', write
$\prod w(\theta_j)=Re^{i\Theta}$ and $\varphi:=\min(\Theta,\pi)$, so that
$\abs{1-Re^{i\varphi}}^2=f_\varphi(R):=1-2R\cos\varphi+R^2$. For $\Theta\le\pi$ this is
the value at $\theta$. For $\Theta>\pi$ it dominates that value, because
$\abs{1-Re^{i\Theta}}\le1+R$.

\emph{Some $\theta_j=\pi/2$.} Then $\alpha=\pi/2$, $R=0$ and the value is $1$, which the
diagonal attains at $t=\pi/2$. Assume from now on that all $\theta_j<\pi/2$.

\emph{Case $\Theta\le\pi$.} First, $R\ge\cos\Theta$: by Lemma~\ref{lem:subadd} if
$\Theta\le\pi/2$, and trivially otherwise. Next, $f_\Theta'(R)=2(R-\cos\Theta)\ge0$ on
$[R,\cos^n(\Theta/n)]$, so Lemma~\ref{lem:jensen} gives
$f_\Theta(R)\le f_\Theta(\cos^n(\Theta/n))=\abs{1-w(\Theta/n)^n}^2$. Finally
$\Theta/n\le\alpha$.

\emph{Case $\Theta>\pi$.} Then $n\alpha>\pi$, so $n\ge3$ and $\pi/n<\alpha$. By
Lemma~\ref{lem:jensen} and the monotonicity of $\cos$,
\[
\abs{1-Re^{i\varphi}}=1+R\le1+\cos^n(\Theta/n)\le1+\cos^n(\pi/n)=\abs{1-w(\pi/n)^n},
\]
since $w(\pi/n)^n=-\cos^n(\pi/n)$.
\end{proof}

\begin{remark}
The split at $\Theta=\pi$ is necessary: for $\Theta>\pi$ the bound $R\ge\cos\Theta$,
and with it the monotonicity in $R$, can fail.
\end{remark}

\section{Upper bound}\label{sec:upper}

\begin{theorem}\label{thm:upper}
$C^P_T(\eta)\le C_k(\eta)$.
\end{theorem}

\begin{proof}
Normalize by (N1) and (N2). If $k=1$, then $R_r=P_rF_r$ and $E^P_T=0$. Otherwise
Corollary~\ref{cor:envelope} and the second part of Lemma~\ref{lem:diagonal}, with
$n=k-1$ and $\alpha=\arcsin\eta$, give $E^P_T\le\eta\,C_k(\eta)$.
\end{proof}

\section{The universal sharp witness}\label{sec:witness}

\begin{theorem}[Universal sharp witness]\label{thm:witness}
Let $T$ be a finite rooted tree with $k$ internal nodes, $\eta\in(0,1]$, and choose
$\theta_v\in[0,\arcsin\eta]$ for the non-root internal nodes. The realization below
is \PMTA-admissible with $M=1$, $\rho=\eta$ and unit leaves, and
\[
E^P_T=\Bigl|1-\prod_{v\in\Vint\setminus\{r\}}w(\theta_v)\Bigr|.
\]
\end{theorem}

\emph{Construction.} All spaces are $\R^2\cong\C$ with $e_0=1$, and all leaves carry
$e_0$. For an internal $v$, let $I_v$ be its internal-child slots and $L_v$ its leaf
slots.
\begin{itemize}
\item At non-root $v$: $P_v(z)=\mathrm{Re}(z)$ (rank one) and
  $\mu_v(x)=e^{i\theta_v}\prod_{j\in I_v}x_j\prod_{\ell\in L_v}\mathrm{Re}(x_\ell)$
  (complex product).
\item At the root: $P_r=I$ and
  $\mu_r(x)=\prod_{j\in I_r}x_j\prod_{\ell\in L_r}\mathrm{Re}(x_\ell)$.
\end{itemize}

\begin{proof}
\emph{Norms.} Since $\abs{\prod x_j}=\prod\abs{x_j}$ and $\abs{\mathrm{Re}(x)}\le\abs x$,
every law has norm $1$.

\emph{Closure at non-root $v$.} Take internal slots in $\Ran P=\R$ and arbitrary leaf
slots. Every factor is real, so the output is $e^{i\theta_v}t$ with $t$ real and
$\abs t\le\prod\norm x$. Hence
$\norm{Q_v\mu_v(x)}=\abs{\sin\theta_v}\abs t\le\eta\prod\norm x$. At the root, $Q_r=0$.

\emph{Evaluation.} By induction, every non-root internal $v$ satisfies
$F_v=e^{i\Phi_v}$ and $R_v=A_v$, where
$\Phi_v=\sum_{u\in\Vint(T_v)}\theta_u$ and $A_v=\prod_{u\in\Vint(T_v)}\cos\theta_u$.
At a leaf-only node, $F_v=e^{i\theta_v}$ and $R_v=\mathrm{Re}(e^{i\theta_v})=\cos\theta_v$.
In the inductive step,
$R_v=\mathrm{Re}\bigl(e^{i\theta_v}\prod_jA_{c_j}\bigr)=\cos\theta_v\prod_jA_{c_j}$
because the $A$'s are real. At the root, $F_r=e^{i\Theta}$ and $R_r=\prod_{v\ne r}\cos\theta_v$
with $\Theta=\sum_{v\ne r}\theta_v$, and $P_r=I$ gives
\[
E^P_T=\bigl|e^{i\Theta}-\textstyle\prod\cos\theta_v\bigr|=\bigl|1-\prod w(\theta_v)\bigr|,
\]
using $\abs{1-\bar z}=\abs{1-z}$.
\end{proof}

\begin{corollary}\label{cor:lower}
$C^P_T(\eta)\ge C_k(\eta)$.
\end{corollary}

\begin{proof}
In Theorem~\ref{thm:witness} take
$\theta_v\equiv\theta_*\in\arg\max_{0\le\theta\le\arcsin\eta}\abs{1-w(\theta)^{k-1}}$.
For $k=1$ both sides are $0$.
\end{proof}

\begin{remark}[The asymptotic witness]
An earlier construction multiplied the phases of leaf slots as well and read an
imaginary part, which gives $\abs{\sin((k-1)\theta)}$. As written, it violates leaf-slot
closure; routing leaves through $\mathrm{Re}(\cdot)$ repairs it. The repaired family proves
$\lim_{\eta\downarrow0}C^P_T(\eta)=k-1$, but for $k\ge3$ it does not attain $C_k(\eta)$.
Indeed,
\[
\abs{1-w(\theta)^n}^2=\sin^2(n\theta)+\bigl(\cos^n\theta-\cos(n\theta)\bigr)^2.
\]
\end{remark}

\section{Proof of Theorem R and scope}\label{sec:proof}

\begin{proof}[Proof of Theorem R]
Theorem~\ref{thm:upper} and Corollary~\ref{cor:lower}.
\end{proof}

\subsection*{Degenerate cases}
\begin{itemize}
\item $\eta=0$ is outside the class.
\item $\eta=1$ with some $\theta=\pi/2$, or some $\rho_i=0$: Lemma~\ref{lem:lemma3} and
  Lemma~\ref{lem:diagonal}.
\item $R_v=0$ or $\tilde R=0$: Lemma~\ref{lem:lemma3}.
\item $P_v\in\{0,I\}$: no division by $\norm{P_vg}$ occurs.
\item $\norm{F_v}<1$: allowed by Lemma~\ref{lem:realization}.
\item Angle labels above $\pi$: the truncation $\min(\cdot,\pi)$ and the case
  $\Theta>\pi$ in Lemma~\ref{lem:diagonal}.
\item Collinear or one-dimensional factors: the reductions in Theorem~\ref{thm:TA}.
\item Arity $0$ and arbitrary arity: Lemma~\ref{lem:lemma3} and Theorem~\ref{thm:TA}.
\item Mixed leaf and internal children: (N2) for the upper bound, and the
  $\mathrm{Re}$-gates for the witness.
\item $k=1$: Theorem~\ref{thm:upper} and Corollary~\ref{cor:lower}.
\end{itemize}

\subsection*{Scope}
\begin{itemize}
\item \emph{Complex spaces.} The upper bound plausibly transfers by realification;
  this is not claimed here. The witness has complex-bilinear norm $\sqrt2$, and the
  complex lower bound is open.
\item \emph{Directed acyclic graphs.} Proposition~\ref{prop:unroll} fails as stated.
\item \emph{Other variants.} Shared laws, fixed ranks and oblique projectors are
  outside \PMTA.
\end{itemize}

\subsection*{Computational controls (not premises)}
\begin{itemize}
\item Random falsification of Lemma~\ref{lem:lemma3}: $344\,000$ node cases at
  arities $1$--$3$, plus adversarial optimization. Every flagged case was a
  discretization artifact of the $\theta$-grid.
\item Cutting-plane LP evaluation of Theorem~\ref{thm:TA} for $m=2,3,4$: the upper
  estimate is within $1.6\cdot10^{-7}$ of the bound.
\item $14\,000$ degenerate node cases at the proof angle: largest violation
  $6.7\cdot10^{-16}$.
\item The witness on $4000$ random trees ($k\le10$, arity $\le4$): attainment error
  $\le3.6\cdot10^{-15}$.
\item The Diagonal Lemma for $n\le12$: no excess beyond grid resolution.
\end{itemize}

\section{Consequences}\label{sec:consequences}

Write $n=k-1$ and $s=\sin^2\theta$. Then $\abs{1-w(\theta)^n}^2=E_k(s)$, where
\[
E_k(s):=1+(1-s)^n-2(1-s)^{n/2}\cos(n\theta)
\]
is a polynomial in $s$, because $\cos^n\theta\cos(n\theta)$ is an even polynomial in
$\cos\theta$. Suppose $E_k$ has a unique critical point $s_c(k)$ in
$(0,\sin^2(\pi/n))$. Then
\[
C_k(\eta)=\begin{cases}\sqrt{E_k(\eta^2)}/\eta,& \eta^2\le s_c(k),\\ \sqrt{E_k(s_c(k))}/\eta,&\eta^2\ge s_c(k).\end{cases}
\]
There are two regimes: below the critical leakage the extremizer uses the whole
budget, and above it the absolute error saturates. Uniqueness of the critical point
was verified in exact arithmetic for $3\le k\le12$.

\begin{center}\small
\renewcommand{\arraystretch}{1.2}\setlength{\tabcolsep}{4pt}
\begin{tabular}{c|l|c|c|c}
$k$ & $E_k(s)$ & $s_c(k)$ & $\max_\eta\eta C_k(\eta)$ & $a_k$\\\hline
$2$ & $s$ & $1$ & $1$ & $0$\\
$3$ & $s(4-3s)$ & $2/3$ & $2/\sqrt3$ & $3/4$\\
$4$ & $s(9-15s+7s^2)$ & $3/7$ & $9/7$ & $5/2$\\
$5$ & $s(4-3s)(4-8s+5s^2)$ & $\tfrac35-\tfrac{\sqrt{21}}{15}$ & $1.38090\ldots$ & $11/2$\\
$6$ & $s(25-100s+160s^2-115s^3+31s^4)$ & $\tfrac{15-\sqrt{70}}{31}$ & $1.45299\ldots$ & $10$\\
$7$ & $s(4-3s)(1-3s+3s^2)(9-15s+7s^2)$ & cubic root & $1.50959\ldots$ & $65/4$
\end{tabular}
\end{center}

\begin{proposition}[Second-order coefficient]\label{prop:ak}
As $\eta\downarrow0$,
\[
C_k(\eta)=(k-1)-a_k\eta^2+O(\eta^4),\qquad a_k=\frac{(k-1)(k-2)(k+6)}{24}.
\]
\end{proposition}

\begin{proof}
For small $\eta$ the maximum is attained at $\theta=\arcsin\eta$, since
$\abs{1-w^n}$ is increasing near $0$. We have
$\log w(\theta)=i\theta-\theta^2/2+O(\theta^4)$, so $X:=n\log w=in\theta-n\theta^2/2+O(\theta^4)$
and
\[
\abs{1-e^X}^2=n^2\theta^2+\bigl(\tfrac{n^2}4-\tfrac{n^3}2-\tfrac{n^4}{12}\bigr)\theta^4+O(\theta^6).
\]
Hence $\abs{1-w^n}=n\theta\bigl(1+(\tfrac18-\tfrac n4-\tfrac{n^2}{24})\theta^2\bigr)+O(\theta^5)$.
Substituting $\theta=\eta+\eta^3/6+O(\eta^5)$ and dividing by $\eta$ gives
$n-\frac{n(n-1)(n+7)}{24}\eta^2+O(\eta^4)$.
\end{proof}

The coefficient grows like $k^3/24$. The sequence $3/4,\,5/2$ at $k=3,4$ does
\emph{not} determine it: the quadratic $(k-2)(2k-3)/4$ fits both values but predicts
$a_5=21/4$ instead of $11/2$.

\begin{proposition}[Large $k$]\label{prop:largek}
Let $G_k:=\max_{\eta\in(0,1]}\eta\,C_k(\eta)$ and let $\eta_c(k)$ be the leakage at
which the maximum over $\theta$ is first attained. Then
$1+\cos^{k-1}\bigl(\tfrac{\pi}{k-1}\bigr)\le G_k<2$, so $G_k\to2$, and
$\eta_c(k)\le\sin\bigl(\tfrac{\pi}{k-1}\bigr)\to0$.
\end{proposition}

\begin{proof}
The proof of Lemma~\ref{lem:diagonal} shows that the maximum of $\abs{1-w^n}$ over
$[0,\pi/2]$ is attained in $[0,\pi/n]$. At $\theta=\pi/n$ we have
$w^n=-\cos^n(\pi/n)$. Moreover $\abs{1-w^n}\le1+\cos^n\theta<2$ for $\theta>0$. Finally
$\cos^n(\pi/n)=\exp(-\pi^2/(2n)+O(n^{-3}))\to1$.
\end{proof}

Numerically, $G_k=2-\pi^2/(2k)+o(1/k)$ and $k\,\eta_c(k)\to\pi$. Neither monotonicity
in $k$ nor these asymptotic constants are proved here.

\begin{remark}[Why the Löwner order O2]
A weaker candidate order compares the rotated distances
$\abs{a-\lambda\zeta e^{i\varphi}}$ and $\abs{1-\lambda ze^{i\varphi}}$ for all
$\lambda\in[0,1]$ and $\varphi\in[0,\pi]$. It contains strictly more states: in random
sampling about $7\%$ of the states it dominates are not O2-dominated. It is not
needed. What makes the induction close is that O2 allows the dilation angle to be
chosen \emph{smaller} than the accumulated label (Lemma~\ref{lem:N-}). An earlier
attempt that forced the angle to equal the label required case-by-case repairs and
did not extend beyond four internal nodes.
\end{remark}

\section*{Questions for the reviewer}

\begin{question}
Is Theorem~\ref{thm:TA} correct for arbitrary finite arity? Is it known, for
instance in the theory of projective tensor norms, contractive dilations, or
scaled relative graphs?
\end{question}

\begin{question}
Does the contraction and dilation argument in Lemma~\ref{lem:lemma3} correctly carry
the raw pair $(f,g)$ to a dominating chain pair at angle $\tilde\psi+\theta_v$? In
particular, is Lemma~\ref{lem:N-} applied under its exact hypotheses?
\end{question}

\begin{question}
Does any step use closure beyond the trajectory form (TC)?
\end{question}

\begin{thebibliography}{9}
\bibitem{RHY} E.~K.~Ryu, R.~Hannah, X.~Yin, \emph{Scaled relative graphs: nonexpansive operators via 2D Euclidean geometry}, Math.\ Program.\ (2022); arXiv:1902.09788.
\bibitem{Ryan} R.~A.~Ryan, \emph{Introduction to Tensor Products of Banach Spaces}, Springer (2002).
\bibitem{NF} B.~Sz.-Nagy, C.~Foia\c{s}, H.~Bercovici, L.~K\'erchy, \emph{Harmonic Analysis of Operators on Hilbert Space}, 2nd ed., Springer (2010).
\end{thebibliography}

\end{document}
